% !TeX root = ../navier-stokes-manuscript.tex
% ============================================================
% 2.1 Vortex Stretching
% ============================================================

The momentum equation has four pieces:
\[
  \underbrace{\partial_t u}_{\text{local acceleration}}
  +
  \underbrace{(u\cdot\nabla)u}_{\text{nonlinear transport}}
  =
  \underbrace{-\nabla p}_{\text{pressure}}
  +
  \underbrace{\nu\Delta u}_{\text{viscous diffusion}}.
\]
Nonlinear transport moves the velocity through the velocity field itself, and in three dimensions that motion can sharpen spatial variation.
Viscosity diffuses the same variation.
The regularity hope is that diffusion always has enough time to smooth whatever transport concentrates.
If that remains true at every scale reached by the fluid, the velocity stays smooth.

\subsection{Vortex Stretching}\label{sub:ThePerfectlyStretchingVortex}

A three-dimensional flow can sharpen through its local rotation.
Write the vorticity as \(\omega:=\nabla\times u\) and the symmetric strain as
\[
  S:=\frac12\bigl(\nabla u+(\nabla u)^{\mathsf T}\bigr),
\]
which records the part of the velocity gradient that changes the shape of a fluid element.
In three dimensions the vorticity can point along a direction in which this strain extends the fluid.

Follow one material segment in such an aligned region.
Let \(e\) be its axial direction and \(L(t)\) its length.
When
\[
  Se=\sigma(t)e,
  \qquad
  \sigma(t)>0,
  \qquad
  \frac{\dot L(t)}{L(t)}
  =
  e\cdot Se
  =
  \sigma(t).
\]
The surrounding flow lengthens the segment along its axis.
Incompressibility fixes what must happen across it at the same time.
Let \(\mathcal V\) be the fixed material volume, set \(A(t):=\mathcal V/L(t)\), and write \(r_{\mathrm{eff}}(t):=\sqrt{A(t)/\pi}\).
Then
\[
  \nabla\cdot u=0,
  \qquad
  \frac{d}{dt}(A L)=0,
  \qquad
  \frac{\dot A}{A}=-\sigma(t),
  \qquad
  \frac{\dot r_{\mathrm{eff}}}{r_{\mathrm{eff}}}
  =
  -\frac{\sigma(t)}{2}.
\]
The segment narrows as it lengthens.

The same alignment acts on the rotation carried inside that contraction.
Taking the curl of the momentum equation gives
\[
  \frac{D\omega}{Dt}
  =
  S\omega+\nu\Delta\omega,
  \qquad
  \frac{D}{Dt}
  :=
  \partial_t+u\cdot\nabla,
\]
where the material derivative carries the rotation with the fluid, viscosity diffuses it, and \(S\omega\) records the change produced by the surrounding strain.
When \(\omega=|\omega|e\), the stretching part of the vorticity equation becomes
\[
  S\omega
  =
  \sigma(t)\omega,
  \qquad
  \frac12\frac{D|\omega|^2}{Dt}
  =
  \sigma(t)|\omega|^2
  +
  \nu\,\omega\cdot\Delta\omega.
\]
The positive term is the amplification produced by axial stretching; viscous diffusion remains present in the same evolution.
On an interval where the full right-hand side stays positive, the rotation strengthens while the cross-section contracts.

This sharpening does not require the segment to carry infinite kinetic energy.
If \(|u|\leq U_0\) on the material segment \(\Omega_{\mathrm{seg}}(t)\), then
\[
  E_{\mathrm{seg}}(t)
  :=
  \frac12\int_{\Omega_{\mathrm{seg}}(t)}|u(x,t)|^2\,dx
  \leq
  \frac12U_0^2\mathcal V
  <
  \infty,
  \qquad
  \left|\frac12\bigl(\nabla u-(\nabla u)^{\mathsf T}\bigr)\right|_{\mathrm F}
  =
  \frac{|\omega|}{\sqrt2}
  \leq
  |\nabla u|_{\mathrm F}.
\]
The velocity can stay bounded on the segment while its rotating gradient grows.
The material volume remains fixed, the width becomes smaller, and finite energy has not removed the concentration.
The original viscosity hope must now be tested at that smaller width, with every term in the equation rescaled together.

\divider{\NavierStokes{} Scaling}

Fix \(t_0\) and set \(\ell:=r_{\mathrm{eff}}(t_0)\).
For \(\lambda>1\), define
\[
  u_\lambda(x,t)
  :=
  \lambda u(\lambda x,\lambda^2t),
  \qquad
  p_\lambda(x,t)
  :=
  \lambda^2p(\lambda x,\lambda^2t).
\]
At time \(\lambda^{-2}t_0\), the vortex in \(u_\lambda\) has width \(\lambda^{-1}\ell\).
This is a comparison solution at a finer scale, not a later state of the material segment just followed.
Its four momentum terms are
\[
\begin{aligned}
  \partial_tu_\lambda(x,t)
  &=
  \lambda^3(\partial_tu)(\lambda x,\lambda^2t),
  \\
  (u_\lambda\cdot\nabla)u_\lambda(x,t)
  &=
  \lambda^3\bigl((u\cdot\nabla)u\bigr)(\lambda x,\lambda^2t),
  \\
  \nabla p_\lambda(x,t)
  &=
  \lambda^3(\nabla p)(\lambda x,\lambda^2t),
  \\
  \nu\Delta u_\lambda(x,t)
  &=
  \lambda^3\nu(\Delta u)(\lambda x,\lambda^2t).
\end{aligned}
\]
At the finer width, viscosity is larger in absolute size, but nonlinear transport is larger by the same \(\lambda^3\) factor.
Pressure and local acceleration scale with them.
Going smaller gives viscosity no automatic advantage.

\divider{Critical Height}

The equation has retained the same balance while the sizes used to measure the solution have changed.
Let \(\Lambda:=(-\Delta)^{1/2}\).
A change of variables gives
\[
  \norm{u_\lambda(t)}_{\dot H^s}^2
  =
  \lambda^{2s-1}
  \norm{u(\lambda^2t)}_{\dot H^s}^2.
\]
The exponent vanishes at \(s=\tfrac12\), selecting the scale-invariant height
\[
  H(t)
  :=
  \frac12\norm{u(t)}_{\dot H^{1/2}}^2
  =
  \frac12\norm{\Lambda^{1/2}u(t)}_2^2.
\]
With \(H_\lambda\) denoting the critical height of \(u_\lambda\),
\[
  \norm{u_\lambda(t)}_2^2
  =
  \lambda^{-1}\norm{u(\lambda^2t)}_2^2,
  \qquad
  H_\lambda(t)=H(\lambda^2t),
  \qquad
  \norm{\nabla u_\lambda(t)}_2^2
  =
  \lambda\norm{\nabla u(\lambda^2t)}_2^2.
\]
At finer scale, kinetic energy falls and the first-derivative norm rises, while the critical height does not move.
Energy can miss the concentration seen by the gradient.
The invariant height is where the competition remains visible without being weakened or enlarged by the change of scale.

\divider{Nonlinear Production versus Viscous Dissipation}

To see that competition in time, apply \(\Lambda^{1/2}\) to the momentum equation and pair with \(\Lambda^{1/2}u\).
The signed nonlinear production is
\[
  \mathcal P_{1/2}(t)
  :=
  -\operatorname{Re}
  \pair
    {\Lambda^{1/2}u(t)}
    {\Lambda^{1/2}\bigl((u(t)\cdot\nabla)u(t)\bigr)}_{L^2}.
\]
The pressure pairing vanishes because \(\nabla\cdot u=0\), while the Laplacian contributes \(-\nu\norm{\Lambda^{3/2}u}_2^2\).
The critical height obeys
\[
  H'(t)
  +
  \nu\norm{\Lambda^{3/2}u(t)}_2^2
  =
  \mathcal P_{1/2}(t).
\]
It rises exactly when nonlinear production exceeds viscous dissipation.
Under the finer-scale comparison, the two contributions become
\[
  \mathcal P_{1/2}[u_\lambda](t)
  =
  \lambda^2\mathcal P_{1/2}[u](\lambda^2t),
  \qquad
  \nu\norm{\Lambda^{3/2}u_\lambda(t)}_2^2
  =
  \lambda^2\nu\norm{\Lambda^{3/2}u(\lambda^2t)}_2^2.
\]
Production and dissipation remain matched in scaling.
Their common \(\lambda^2\) factor is the inverse of the time contraction \(t\mapsto\lambda^{-2}t\): the finer balance runs on a shorter clock.

\divider{The Heat Clock}

Viscosity supplies that width-dependent clock through the linear heat equation \(v_t=\nu\Delta v\):
\[
  \widehat v(\xi,t+\delta t)
  =
  e^{-\nu|\xi|^2\delta t}\widehat v(\xi,t).
\]
A scale-localized structure of width \(r\) is concentrated around frequencies \(|\xi|\simeq r^{-1}\).
Its viscous change becomes order one when \(\nu|\xi|^2\delta t\simeq1\), giving the comparison time
\[
  \tau_\nu(r)
  :=
  \frac{r^2}{\nu}.
\]
This clock measures how quickly diffusion acts at width \(r\); it does not assert that the nonlinear flow follows the heat equation.
At the finer width,
\[
  \tau_\nu(\lambda^{-1}r)
  =
  \lambda^{-2}\tau_\nu(r),
\]
the same time contraction carried by the full \NavierStokes{} scaling.

\divider{The Zeno Cascade}

Each further contraction of width brings a further contraction of the comparison time.
For the geometric sequence
\[
  r_n:=2^{-n}r_0,
  \qquad
  \tau_n:=\frac{r_n^2}{\nu},
\]
the heat clocks satisfy
\[
  \tau_n
  =
  \frac{r_0^2}{\nu}4^{-n},
  \qquad
  \sum_{n=0}^{\infty}\tau_n
  =
  \frac{4r_0^2}{3\nu}
  <
  \infty.
\]
The sum shows only that these clocks do not prevent infinitely many scale transitions from fitting into finite time; it does not construct such an evolution.
For the sequence to describe a candidate singularity, one \NavierStokes{} solution would have to pass through widths
\[
  r_0>r_1>r_2>\cdots\longrightarrow0
\]
at times
\[
  t_0<t_1<t_2<\cdots\uparrow T_*<\infty,
  \qquad
  t_{n+1}-t_n\asymp\tau_n,
\]
with comparison constants independent of \(n\).
The remaining time would then satisfy
\[
  T_*-t_n
  \asymp
  \sum_{k=n}^{\infty}\tau_k
  =
  \frac{4r_n^2}{3\nu}.
\]
Under that candidate schedule, each finer width must be formed during an interval tending to zero at order \(\tau_n\).
Passing from one width to the next changes \(u\) over that shrinking interval; the rate of change, the rate of that rate, and every later response now belong to the same demand.
The spatial sequence has become a demand on \(u,\partial_tu,\partial_t^2u,\partial_t^3u,\ldots\): velocity, acceleration, jerk, and the temporal Tower.
